\section{Results} \label{result}

\subsection{Simulation of the Quantum Component}

This section presents simulation results of the BB84 quantum component with increasing levels of complexity.

\subsubsection{Local Simulation: Noiseless and without Eavesdropping}

The BB84 QKD simulation is executed locally five times using 20 qubits. The outputs presented below correspond to one representative run.

\noindent
\resizebox{\columnwidth}{!}{%
$\begin{array}{l l}
\textbf{Alice's Bits:}       & 0,1,1,0,1,0,0,0,0,0,0,1,0,1,0,0,1,1,1,1 \\
\textbf{Alice's Bases:}      & Z,X,Z,Z,Z,X,X,Z,Z,X,Z,X,Z,X,X,Z,X,X,Z,X \\
\textbf{Bob's Bases:}        & X,Z,Z,Z,Z,Z,Z,X,Z,Z,Z,Z,Z,Z,Z,X,Z,X,Z,X \\
\textbf{Bob's Bits:}         & 0,0,1,0,1,0,1,1,0,1,0,1,0,1,1,0,0,1,1,1 \\
\textbf{Alice's Sifted Key}  & 1, 0, 1, 0, 0, 0, 1, 1, 1 \\
\textbf{Bob's Sifted Key}    & 1, 0, 1, 0, 0, 0, 1, 1, 1 \\
\end{array}$%
}

Since Alice and Bob's sifted keys are identical, QBER $= 0$. 

\subsubsection{Local Simulation: Noisy but without Eavesdropping}

We add depolarizing noise to the circuit that simulates Alice's preparations and Bob's measurements.
We run this simulation five times and observe the QBER. The outputs presented below are for one representative run.

\noindent
\resizebox{\columnwidth}{!}{%
$\begin{array}{l l}
\textbf{Alice's Bits:}       & 0, 1, 1, 0, 1, 0, 0, 0, 0, 0, 0, 1, 0, 1, 0, 0, 1, 1, 1, 1 \\
\textbf{Alice's Bases:}      & Z, X, Z, Z, Z, X, X, Z, Z, X, Z, X, Z, X, X, Z, X, X, Z, X \\
\textbf{Bob's Bases:}        & X, Z, Z, Z, Z, Z, Z, X, Z, Z, Z, Z, Z, Z, Z, X, Z, X, Z, X \\
\textbf{Bob's Bits:}         & 1, 0, 1, 1, 1, 1, 0, 0, 0, 0, 0, 1, 0, 0, 0, 1, 1, 1, 0, 1 \\
\textbf{Alice's Sifted Key}  & 1, 0, 1, 0, 0, 0, 1, 1, 1 \\
\textbf{Bob's Sifted Key}    & 1, 1, 1, 0, 0, 0, 1, 0, 1 \\
\end{array}$%
}

Since Alice and Bob’s sifted keys differ in 2 out of 9 bits, the QBER is $2/9 \approx 22\%$. This increase reflects the added noise. 
Because this exceeds the 11\% threshold, Alice and Bob would abort the protocol in this case.

\subsubsection{IBM Hardware Simulation: Noisy but without Eavesdropping}

This simulation was run on an IBM Quantum Platform repeatedly 10 times. The outputs presented below correspond to one representative run.

\noindent
\resizebox{\columnwidth}{!}{%
$\begin{array}{l l}
\textbf{Alice's Bits:}       & 0, 1, 1, 0, 1, 0, 0, 0, 0, 0, 0, 1, 0, 1, 0, 0, 1, 1, 1, 1 \\
\textbf{Alice's Bases:}      & Z, X, Z, Z, Z, X, X, Z, Z, X, Z, X, Z, X, X, Z, X, X, Z, X \\
\textbf{Bob's Bases:}        & X, Z, Z, Z, Z, Z, Z, X, Z, Z, Z, Z, Z, Z, Z, X, Z, X, Z, X \\
\textbf{Bob's Bits:}         & 0, 0, 0, 0, 1, 0, 0, 0, 0, 0, 0, 0, 0, 0, 1, 1, 1, 1, 1, 1 \\
\textbf{Alice's Sifted Key}  & 1, 0, 1, 0, 0, 0, 1, 1, 1 \\
\textbf{Bob's Sifted Key}    & 0, 0, 1, 0, 0, 0, 1, 1, 1 \\
\end{array}$%
}

Alice and Bob's sifted keys differ in one bit, which amounts to a QBER of approximately 11\%. 
However, QBER $= 0$ in the remaining nine runs. Hence, this IBM machine is almost noiseless for tasks like preparations and measurements.

\subsubsection{Local Simulation: Eavesdropping but Noiseless}

We now show how Eve's intercept-resend attack can result in a much higher QBER. Here, we run the simulation locally five times. 
Below are the outputs corresponding to one representative run.

\noindent
\resizebox{\columnwidth}{!}{%
$\begin{array}{l l}
\textbf{Alice's Bits:}       & 0, 1, 1, 0, 1, 0, 0, 0, 0, 0, 0, 1, 0, 1, 0, 0, 1, 1, 1, 1 \\
\textbf{Alice's Bases:}      & Z, X, Z, Z, Z, X, X, Z, Z, X, Z, X, Z, X, X, Z, X, X, Z, X \\
\textbf{Eve's Bases:}        & X, Z, Z, Z, Z, Z, Z, X, Z, Z, Z, Z, Z, Z, Z, X, Z, X, Z, X \\
\textbf{Eve's Bits:}         & 0, 1, 1, 0, 1, 0, 0, 0, 0, 0, 0, 0, 0, 0, 0, 0, 1, 1, 1, 1 \\
\textbf{Bob's Bases:}        & X, X, X, X, Z, Z, Z, X, Z, X, X, Z, Z, X, Z, X, X, Z, X, X \\
\textbf{Bob's Bits:}         & 0, 0, 0, 1, 1, 0, 0, 0, 0, 0, 1, 0, 0, 0, 0, 0, 1, 1, 0, 1 \\
\textbf{Alice's Sifted Key}  & 1, 1, 0, 0, 0, 1, 1, 1\\
\textbf{Bob's Sifted Key}    & 0, 1, 0, 0, 0, 0, 1, 1 \\
\end{array}$%
}

The QBER rises to 25\%, which coincides with the error rate discussed in section II.C. In addition, the average QBER over all five runs is approximately 25\%.


\subsubsection{IBM Hardware Simulation: Eavesdropping with Noise}

Finally, we simulate the intercept-resend attack by executing the BB84 protocol five times on IBM Quantum Platform.
Below are the outputs for one representative run.

\noindent
\resizebox{\columnwidth}{!}{%
$\begin{array}{l l}
\textbf{Alice's Bits:}       & 0, 1, 1, 0, 1, 0, 0, 0, 0, 0, 0, 1, 0, 1, 0, 0, 1, 1, 1, 1 \\
\textbf{Alice's Bases:}      & Z, X, Z, Z, Z, X, X, Z, Z, X, Z, X, Z, X, X, Z, X, X, Z, X \\
\textbf{Eve's Bases:}        & X, Z, Z, Z, Z, Z, Z, X, Z, Z, Z, Z, Z, Z, Z, X, Z, X, Z, Z \\
\textbf{Eve's Bits:}         & 1, 1, 1, 0, 1, 0, 0, 0, 0, 1, 0, 0, 0, 1, 0, 0, 1, 1, 1, 1 \\
\textbf{Bob's Bases:}        & X, X, X, X, Z, Z, Z, X, Z, X, X, Z, Z, X, Z, X, X, Z, X, X \\
\textbf{Bob's Bits:}         & 1, 1, 0, 1, 1, 0, 0, 0, 0, 1, 1, 0, 0, 0, 0, 0, 0, 1, 1, 1 \\
\textbf{Alice's Sifted Key}  & 1, 1, 0, 0, 0, 1, 1, 1 \\
\textbf{Bob's Sifted Key}    & 1, 1, 0, 1, 0, 0, 0, 1 \\
\end{array}$%
}

As in the preceding local simulation case, the QBER for this run is 37.5\%, while the average over all five runs is 27.5\%.

\subsection{Simulation of the Classical Post-Processing Component}

This section shows the results of two key classical post-processing tasks: error correction and privacy amplification. 
We also demonstrate one-time pad cryptography by showing how the plaintext "BB84" can be encrypted by Alice and decrypted by Bob using the shared key.

\subsubsection{Error Correction: A Cascade Protocol}

The main inputs to the function cascade\_protocol are Alice and Bob's sifted keys. Using the above simulation of BB84 with sufficiently low noise, we generated:

\textbf{Alice's sifted key:}  [1, 1, 1, 1, 0, 0, 0, 0, 0, 0, 1, 1, 0, 1, 1, 1, 0, 0, 1, 1, 0, 1, 1, 1, 1, 0, 0, 0, 1, 1, 0, 1, 0, 0, 1, 1, 1, 1, 1, 0, 0, 1, 0, 0, 1, 1, 0, 0, 1, 0, 0, 0, 1, 1, 0, 1, 1, 1, 1, 0, 0, 1, 0, 1]

\textbf{Bob's sifted key:}  [1, 1, 1, 1, 0, 0, 0, 0, 0, 0, 1, 1, 0, 1, 1, 1, 0, 0, 1, 1, 1, 1, 1, 1, 1, 0, 0, 0, 1, 1, 0, 1, 0, 0, 1, 1, 0, 1, 1, 0, 0, 1, 0, 0, 1, 1, 0, 0, 1, 0, 0, 0, 1, 1, 0, 0, 0, 1, 1, 0, 0, 0, 0, 1].

The Cascade protocol first estimates QBER using a few bits selected from Alice and Bob's sifted keys. This requires Bob to give up a small number of his bits by sending them to Alice to estimate QBER. 
Here, we selected the first 10 bits. If QBER $< 11\%$, error correction is performed on the sifted key.
Cascade makes multiple passes with different block sizes in a way that errors corrected on the $i^{th}$ pass can reveal the location of more errors in another pass $i'$ where $i\ne i'$. 
Cascade locates five mismatches, which are called $leakage$ because the non-zero parity check bits may reveal information to Eve, and corrects them. 

\textbf{Bob's corrected key:} [1, 1, 1, 1, 0, 0, 0, 0, 0, 0, 1, 1, 0, 1, 1, 1, 0, 0, 1, 1, 0, 1, 1, 1, 1, 0, 0, 0, 1, 1, 0, 1, 0, 0, 1, 1, 1, 1, 1, 0, 0, 1, 0, 0, 1, 1, 0, 0, 1, 0, 0, 0, 1, 1, 0, 1, 1, 1, 1, 0, 0, 1, 0, 1],

\noindent which matches Alice's sifted key.

\subsubsection{Privacy Amplification: A Random Binary Toeplitz Matrix}

The main task is to construct a random binary Toeplitz matrix of an appropriate size. First, we remove the 10 bits from the shared key that were used to estimate QBER during error correction. 
Let $N$ be the key length. The number of rows is $N-leakage-s$, where $leakage=5$ was determined during error correction and $s=10$ is the security margin. 
Thus, our Toeplitz matrix is of size $(N-15)\times N$. The final shared key is the product of this matrix with the $N$-bit shared key.

\textbf{Final shared key}: [1, 0, 0, 0, 1, 1, 0, 1, 1, 1, 1, 0, 1, 1, 1, 0, 0, 1, 1, 0, 1, 0, 1, 1, 1, 1, 1, 0, 1, 1, 0, 1, 0, 1, 1, 1, 1, 0, 0].

\subsubsection{One-Time Pad Encryption, Decryption}

To encrypt the message "BB84", Alice truncates the final shared key to 32 bits and XORs it with the binary representation of "BB84". 

\textbf{Ciphertext in bits}: 11001111101011000101001111011001.

To decrypt the ciphertext, Bob truncates the final shared key and XORs it with the ciphertext. 
We have assumed that Bob knows the plaintext's length. In practice, Alice pads the plaintext with 0's before encryption for Bob to extract the original message from the decrypted bit string.

